\documentclass[11pt]{article}
\usepackage[a4paper,margin=25mm]{geometry}
\usepackage{amsmath,amssymb,amsthm}
\usepackage{longtable,booktabs,array}
\usepackage[hidelinks]{hyperref}
\hypersetup{pdftitle={E96191221685552745\_2: 8x\textasciicircum 4+x+y\textasciicircum 3+y-1=0 has no integer solutions}}
\newtheorem{theorem}{Theorem}
\newtheorem{lemma}[theorem]{Lemma}
% keep the space around binary operators in inline formulas
\medmuskip=4mu plus 2mu\relax
\emergencystretch=1.5em
\tolerance=400
\allowdisplaybreaks
\newcommand{\Z}{\mathbb{Z}}
\newcommand{\Q}{\mathbb{Q}}
\newcommand{\F}{\mathbb{F}}
% Legendre and Jacobi symbols: one fixed size in running text, one in displays
\newcommand{\leg}[2]{\mathchoice{\biggl(\dfrac{#1}{#2}\biggr)}{(#1/#2)}{(#1/#2)}{(#1/#2)}}
\title{\textbf{E96191221685552745\_2}\\[4pt] {\Large The equation $8x^{4}+x+y^{3}+y-1=0$ has no integer solutions}}
\author{}
\date{}
\begin{document}
\maketitle
\vspace{-8ex}
\begin{center}
Size $h=141$, length $l=12$
\end{center}

\begin{theorem}\label{thm:main}
The equation
\begin{equation}\label{eq:main}
8x^{4}+x+y^{3}+y-1=0
\end{equation}
has no integer solutions.
\end{theorem}

The proof uses the cubic Hilbert reciprocity law over the field $K=\Q(\omega)$, where $\omega^2+\omega+1=0$. To a hypothetical integer solution we attach the values of 25 auxiliary polynomials $H_0,\dots,H_{24}$ with coefficients in $\Z[\omega]$, constants $c_0,\dots,c_{24}\in\Z[\omega]$, and a sum $B$ of cubic Hilbert symbols of these values, recorded by a graph with 70 weighted edges. By the reciprocity law, the local contributions $B_p$ of all primes $p$ add up to $0$ modulo $3$. We show that $B_p=0$ for $p\notin S$, compute $B_p$ for the primes $p\in S$, and obtain a contradiction. Two standard facts from class field theory are quoted without proof: the explicit formula for the cubic Hilbert symbol at the places not above $3$, and the Hilbert reciprocity law. Apart from these, the proof uses polynomial identities, which can be checked by expanding both sides, and finite computations with residues.


\section{\texorpdfstring{Cubic Hilbert symbols over $\Q(\omega)$}{Cubic Hilbert symbols over Q(w)}}\label{sec:symbols}

\paragraph{The field.} Let $K=\Q(\omega)$, where $\omega^2+\omega+1=0$. Its ring of integers is $\Z[\omega]$, and every element of $K$ can be written uniquely as $A+B\omega$ with $A,B\in\Q$. We write
\[
\langle A,B\rangle=A+B\omega,
\]
also when $A$ and $B$ are polynomials in $x$ and $y$. Then $\langle A,B\rangle\langle C,D\rangle=\langle AC-BD,\ AD+BC-BD\rangle$, and the norm is $N(\langle A,B\rangle)=A^2-AB+B^2$. In particular, a value $\langle A,B\rangle$ with $A,B\in\Z$ is zero if and only if $A=B=0$.

\paragraph{Places.} The finite places of $K$ correspond to the non-zero prime ideals of $\Z[\omega]$. A prime $p\equiv2\pmod3$ is inert: there is one place above $p$, with uniformizer $p$ and residue field $\Z[\omega]/p\Z[\omega]$ with $p^2$ elements. If $p\equiv1\pmod3$, the polynomial $t^2+t+1$ has two roots $r$ modulo $p$, and for each of them there is a place $v_r$ above $p$, at which $\omega$ is the root of $t^2+t+1$ in $\Z_p$ congruent to $r$ modulo $p$; the completion is $\Q_p$, the uniformizer is $p$, the residue field is $\F_p$, and $\omega$ reduces to $r$. The prime $3$ is ramified: $\pi=1-\omega$ generates the unique prime ideal above $3$, $\pi^2=-3\omega$, and $\Z[\omega]/\pi\Z[\omega]=\F_3$, where $\langle A,B\rangle$ reduces to $A+B$. The only infinite place of $K$ is complex. For a finite place $v$ we write $K_v$ for the completion and $v(a)$ for the normalized valuation, and we call $a$ a \emph{unit} at $v$ if $v(a)=0$.

\paragraph{Hilbert symbols.} For a place $v$ of $K$ and $a,b\in K_v^\times$, the cubic Hilbert symbol $(a,b)_v$ is a cube root of unity, see \cite[Chapter~V, \S3]{N}. We write $(a,b)_v=\omega^{\beta_v(a,b)}$ with $\beta_v(a,b)\in\Z/3\Z$, and we use the following three facts \cite[Chapter~V, \S3 and Chapter~VI, \S8]{N}.
\begin{itemize}
\item[(H1)] $\beta_v(a,bc)=\beta_v(a,b)+\beta_v(a,c)$, $\beta_v(a,b)=-\beta_v(b,a)$, and $\beta_v(a,c^3)=0$. In particular, $\beta_v(a,b)$ depends only on the classes of $a$ and $b$ in $K_v^\times/K_v^{\times3}$.
\item[(H2)] Let $v$ be a finite place not above $3$, let $q$ be the number of elements of its residue field, and let $\varpi$ be a uniformizer at $v$. Write $a=\varpi^eu$ and $b=\varpi^fw$ with units $u$ and $w$, and define $\chi(u)\in\Z/3\Z$ by $\bar u^{(q-1)/3}=\bar\omega^{\chi(u)}$, where the bar denotes the reduction to the residue field. Then $\beta_v(a,b)=e\,\chi(w)-f\,\chi(u)$. In particular, $\beta_v(a,b)=0$ if $a$ and $b$ are both units.
\item[(H3)] (Hilbert reciprocity.) For $a,b\in K^\times$, $\beta_v(a,b)=0$ for all but finitely many places $v$, and $\sum_v\beta_v(a,b)=0$, where the sum is over all places of $K$. At the complex place, $\beta_v(a,b)=0$.
\end{itemize}
The general formula in (H2) contains an additional term $ef\chi(-1)$, which vanishes because $-1=(-1)^3$ is a cube. Formula (H2) fixes the normalization of the symbols; with the opposite sign convention all values $\beta_v$ change sign, and the proof below is the same.

\paragraph{Classes at the places not above $3$.} Let $p\neq3$ be a prime and $v$ a place above $p$. For $a=p^eu\in K_v^\times$ with a unit $u$, put $c_v(a)=(e\bmod3,\ \chi(u))\in\F_3^2$. By (H2) and (H1), $c_v$ is additive, $c_v(a)$ determines the class of $a$ modulo cubes, and $\beta_v(a,b)=e\chi(w)-f\chi(u)$ for $c_v(a)=(e,\chi(u))$ and $c_v(b)=(f,\chi(w))$. We put $c_p(a)=c_v(a)\in\F_3^2$ if $p\equiv2\pmod3$, and $c_p(a)=(c_{v_r}(a),c_{v_{r'}}(a))\in\F_3^4$ if $p\equiv1\pmod3$, where $r<r'$ are the roots of $t^2+t+1$ in $\{1,\dots,p-1\}$; accordingly, $\beta_p$ is the sum of the forms $\beta_v$ over the places $v$ above $p$. For $p\equiv2\pmod3$ and a rational integer $u$ prime to $p$, we have $\bar u\in\F_p$, so that $\bar u^{(p^2-1)/3}=(\bar u^{p-1})^{(p+1)/3}=1$ and $\chi(u)=0$.

\paragraph{The place above $3$.}
\begin{lemma}\label{lem:M3}
Let $v$ be the place above $3$. Every non-zero element of $K_v$ is equal to $\pi^a\omega^b2^c(1+3\omega)^d$ times a cube, for unique $a,b,c,d\in\Z/3\Z$, and every unit congruent to $1$ modulo $9$ is a cube. Put $c_3(A)=(a,b,c,d)\in\F_3^4$. Then $c_3$ is additive, and $\beta_v(A,B)=c_3(A)^{\mathsf T}M_3\,c_3(B)$, where
\[
M_3=\begin{pmatrix}0&0&2&0\\0&0&1&2\\1&2&0&0\\0&1&0&0\end{pmatrix}\pmod3.
\]
\end{lemma}

\begin{proof}
For $\alpha\in\Z[\omega]$, the polynomial $g(s)=s+3s^2+3s^3-\alpha$ satisfies $g'(s)\equiv1\pmod\pi$, so by Hensel's lemma it has a root $s$ in the completion of $\Z[\omega]$, and then $(1+3s)^3=1+9(s+3s^2+3s^3)=1+9\alpha$. Hence every unit congruent to $1$ modulo $9$ is a cube in $K_v$, and the class of a unit modulo cubes depends only on its residue modulo $9$. There are $54$ units in $\Z[\omega]/(9)$, and a direct enumeration shows that the $27$ products $\omega^b2^c(1+3\omega)^d$ with $b,c,d\in\{0,1,2\}$ are pairwise distinct modulo $9$, also after multiplication by $-1=(-1)^3$. Hence they represent all $27$ classes of units modulo cubes. Since $\pi^3$ is a cube times a unit, the exponent of $\pi$ is the valuation modulo $3$, and the first statement follows.

By (H1), it remains to compute $\beta_v$ on the $16$ ordered pairs of the elements $\pi$, $\omega$, $2$ and $1+3\omega$. Their norms are $3$, $1$, $4$ and $7$, hence, for each such pair, all the symbols at the places not above $2$, $3$ and $7$ vanish by (H2), and, by (H3), $\beta_v$ at the place above $3$ is minus the sum of the values at the place above $2$ and at the two places above $7$. These places correspond to $\bar\omega=2$ and $\bar\omega=4$ in $\F_7$. The classes $(e,\chi)$ of the four elements at these places are
\[
\begin{array}{c|ccc}
 & 2 & \bar\omega=2 & \bar\omega=4\\\hline
\pi & (0,2) & (0,0) & (0,2)\\
\omega & (0,1) & (0,2) & (0,2)\\
2 & (1,0) & (0,2) & (0,1)\\
1+3\omega & (0,2) & (1,0) & (0,0)
\end{array}
\]
where, at the places above $7$, $\chi(u)$ is defined by $\bar u^2=\bar\omega^{\chi(u)}$ in $\F_7$, and $7$ is a uniformizer. Formula (H2) and the sum over these three places give $M_3$.
\end{proof}

We write $\beta_3$ for the form of Lemma~\ref{lem:M3}. Since $3=-\omega^2\pi^2$ and $4=2^2$, we have $c_3(3)=(2,2,0,0)$ and $c_3(4)=(0,0,2,0)$.

\paragraph{Residue classes.} The local computations use the following description of the classes of the values of a polynomial on a residue class.

\begin{lemma}\label{lem:residue}
Let $p$ be a prime, $k\geq1$, and let $z\in\Z[\omega]$ be non-zero with $z\equiv\zeta\pmod{p^k}$, where $\zeta=\langle A,B\rangle$ with $0\leq A,B<p^k$. Then $c_p(z)\in b+D$ for the vector $b$ and the subspace $D$ given as follows.
\begin{itemize}
\item[(i)] $p\equiv2\pmod3$. If $\zeta\neq0$, then $b=c_p(\zeta)$ and $D=0$. If $\zeta=0$, then $b=0$, and $D$ is $\F_3^2$, or the span of $(1,0)$ if $z\in\Z$.
\item[(ii)] $p\equiv1\pmod3$. For each root $r$, let $\rho$ be the root of $t^2+t+1$ modulo $p^k$ congruent to $r$. If $A+B\rho\not\equiv0\pmod{p^k}$, the two coordinates of $c_{v_r}(z)$ are those of $c_{v_r}(\zeta)$; otherwise these two coordinates are arbitrary.
\item[(iii)] $p=3$, $z\notin\Z$. If $\zeta\neq0$, let $a=v(\zeta)<2k$ be its valuation at $\pi$ and $m=2k-a$. Then $b=c_3(\zeta)$, and $D=0$ if $m\geq4$, while $D=D_m$ if $m\leq3$, where $D_1$ is spanned by $(0,1,0,0)$, $(0,0,1,0)$ and $(0,0,0,1)$, $D_2$ by $(0,0,1,0)$ and $(0,0,0,1)$, and $D_3$ by $(0,0,1,1)$. If $\zeta=0$, then $b=0$ and $D=\F_3^4$.
\item[(iv)] $p=3$, $z\in\Z$, so that $B=0$. If $3^k\nmid A$, let $j=v_3(A)$. Then $b=c_3(A)$ and $D=0$ if $k-j\geq2$, while $b=j\,c_3(3)$ and $D$ is spanned by $c_3(4)$ if $k-j=1$. If $3^k\mid A$, then $b=0$ and $D$ is spanned by $c_3(3)$ and $c_3(4)$.
\end{itemize}
\end{lemma}

\begin{proof}
(i) If $\zeta\neq0$, then $e=v(z)=v(\zeta)<k$ and $z/p^e\equiv\zeta/p^e\pmod p$, so that $\chi$ is determined. If $z\in\Z$, its unit part is a rational integer, and $\chi=0$. (ii) At $v_r$, $z$ and $\zeta$ are the $p$-adic numbers $A'+B'\omega$ and $A+B\omega$ with $\omega\equiv\rho\pmod{p^k}$, and the argument of (i) applies. (iii) If $\zeta\neq0$, then $v(z)=a$, and $z/\pi^a\equiv\zeta/\pi^a\pmod{\pi^m}$. By Lemma~\ref{lem:M3}, the class of a unit depends only on its residue modulo $9=\omega\pi^4$; the units congruent to $1$ modulo $\pi^m$ form a subgroup, and a direct enumeration of their residues modulo $9$ shows that their classes form $D_m$. (iv) If $3^k\nmid A$, then $z=3^ju$ with a rational unit $u\equiv A/3^j\pmod{3^{k-j}}$. The class of a rational unit depends only on its residue modulo $9$, the units $\pm1$ are cubes, and the rational units congruent to $1$ modulo $3$ are $4^s$ times a unit congruent to $1$ modulo $9$; hence the classes are as stated, and if $3^k\mid A$ the valuation $j$ is also unknown.
\end{proof}

\begin{lemma}\label{lem:envelope}
Let $p$ be a prime, and consider vertices $i$ with vectors $b_i$, subspaces $D_i$ and vectors $\gamma_i$, and a set of edges $(i,j)$ with $i<j$ and weights $w_{ij}\in\{1,2\}$. Put $\rho_{ij}=w_{ij}$ and $\rho_{ji}=-w_{ij}$ for the edges, and $\sigma_i=\gamma_i+\sum_j\rho_{ij}b_j$. Suppose that $\beta_p(d,\sigma_i)=0$ for every $i$ and every $d\in D_i$, and that $\beta_p(d,d')=0$ for every edge $(i,j)$ and all $d\in D_i$, $d'\in D_j$. Then
\[
\sum_{(i,j)}w_{ij}\beta_p(c_i,c_j)+\sum_i\beta_p(c_i,\gamma_i)=\sum_{(i,j)}w_{ij}\beta_p(b_i,b_j)+\sum_i\beta_p(b_i,\gamma_i)
\]
whenever $c_i\in b_i+D_i$ for every $i$. By bilinearity, it suffices to check the hypotheses for $d$ and $d'$ in spanning sets of $D_i$ and $D_j$.
\end{lemma}

\begin{proof}
Write $c_i=b_i+d_i$ with $d_i\in D_i$. As $\beta_p$ is bilinear and $\beta_p(u,v)=-\beta_p(v,u)$, the left side equals the right side plus $\sum_i\beta_p(d_i,\sigma_i)+\sum_{(i,j)}w_{ij}\beta_p(d_i,d_j)$, and both sums vanish.
\end{proof}
\section{The graph}

Put $F(x,y)=8x^{4}+x+y^{3}+y-1$, so that \eqref{eq:main} is the equation $F(x,y)=0$. The certificate consists of 25 polynomials $H_i$ with coefficients in $\Z[\omega]$ and constants $c_i\in\Z[\omega]$, listed in Table~\ref{tab:vert1}, and of 70 \emph{edges} $(i,j)$ with $i<j$ and weights $w_{ij}\in\{1,2\}$, listed below. For an edge $(i,j)$ we put $\rho_{ij}=w_{ij}$ and $\rho_{ji}=-w_{ij}\bmod 3$, and $\rho_{ij}=0$ if $i$ and $j$ are not joined by an edge.
The edges $(i,j;w_{ij})$ are $(0,4;1)$, $(0,10;1)$, $(0,11;2)$, $(0,15;2)$, $(0,16;1)$, $(1,4;1)$, $(1,5;2)$, $(1,6;1)$, $(1,15;2)$, $(1,19;1)$, $(1,24;2)$, $(2,6;2)$, $(2,10;2)$, $(2,12;1)$, $(2,19;1)$, $(2,21;2)$, $(3,13;2)$, $(3,16;1)$, $(3,18;1)$, $(3,21;2)$, $(4,15;1)$, $(4,17;2)$, $(4,18;1)$, $(4,20;2)$, $(4,21;2)$, $(5,11;1)$, $(5,12;1)$, $(5,21;2)$, $(6,10;2)$, $(6,11;1)$, $(6,13;2)$, $(6,14;1)$, $(6,16;1)$, $(6,19;1)$, $(6,20;1)$, $(7,11;2)$, $(7,12;2)$, $(7,16;1)$, $(7,20;2)$, $(8,11;2)$, $(8,14;2)$, $(8,17;2)$, $(8,19;1)$, $(9,12;1)$, $(9,14;1)$, $(9,15;2)$, $(9,17;1)$, $(10,20;2)$, $(10,22;1)$, $(10,23;2)$, $(11,12;1)$, $(12,15;1)$, $(13,14;1)$, $(13,19;1)$, $(13,24;1)$, $(14,16;2)$, $(14,17;1)$, $(14,19;2)$, $(14,22;2)$, $(15,16;1)$, $(15,17;2)$, $(15,18;1)$, $(15,20;1)$, $(17,18;1)$, $(17,21;1)$, $(17,23;1)$, $(18,24;2)$, $(19,20;2)$, $(19,21;2)$, $(20,24;1)$.
Let $(\xi,\eta)$ denote $(x,y)$ or $(y,x)$. We call a vertex $i$ \emph{linear} if $H_i=\varepsilon_i(\lambda_i\eta-s_i(\xi))$ for a unit $\varepsilon_i$ of $\Z[\omega]$, an element $\lambda_i\in\Z[\omega]$ whose norm is a power of $3$, and a polynomial $s_i$ with coefficients in $\Z[\omega]$; the last column of Table~\ref{tab:vert1} gives $\lambda_i\eta=s_i(\xi)$ for the linear vertices, where $\lambda_i=1$ if no factor is shown. For a linear vertex $i$, we put $f_i(t)=\lambda_i^{d}F$ and $g_{ij}(t)=\lambda_i^{d_j}H_j$ evaluated at $\xi=t$, $\eta=s_i(t)/\lambda_i$, where $d$ and $d_j$ are the degrees of $F$ and $H_j$ in $\eta$; these are polynomials in $t$ with coefficients in $\Z[\omega]$.
For every linear vertex $i$, Table~\ref{tab:star} gives an integer $D_i\geq1$ and a polynomial $w_i(t)$ with coefficients in $\Z[\omega]$ such that
\begin{equation}\label{eq:starlin}
D_i^3\,\lambda_i^{e_i}c_i\prod_jg_{ij}^{\rho_{ij}}-w_i^3=f_iQ_i
\end{equation}
for a polynomial $Q_i(t)$ with coefficients in $\Z[\omega]$, where $\rho_{ij}\in\{0,1,2\}$; the table also gives an integer $e_i\geq0$ with $e_i+\sum_j\rho_{ij}d_j\equiv0\pmod3$, and $e_i=0$ if $\lambda_i=1$; $Q_i$ is the quotient of the division of the left side by $f_i$, and the remainder is zero.
For every other vertex $i$, Table~\ref{tab:star} gives an integer $D_i\geq1$ and polynomials $W_i$, $A_i$ and $B_i$ in $x$ and $y$ with coefficients in $\Z[\omega]$ such that
\begin{equation}\label{eq:star}
D_i^3\,c_i\prod_jH_j^{\rho_{ij}}-W_i^3=A_iF+B_iH_i.
\end{equation}
For the edges $(i,j)$ listed in Table~\ref{tab:sep}, the table gives a linear endpoint $k\in\{i,j\}$ and the factorization of the norm of the resultant $R_{ij}=\operatorname{Res}_t(f_k,g_{kl})$, where $\{k,l\}=\{i,j\}$; it is non-zero.
These identities can be checked by expanding both sides and by polynomial division, using $\omega^2=-1-\omega$.

\section{Proof of Theorem~\ref{thm:main}}

Suppose that $(x,y)\in\Z^2$ is a solution of \eqref{eq:main}, so that $F(x,y)=0$; from now on the polynomials denote their values at $(x,y)$, which lie in $\Z[\omega]$.
\paragraph{Step 1: the values are non-zero.} For $i\in\{0,\allowbreak 1,\allowbreak 2,\allowbreak 4,\allowbreak 5,\allowbreak 6,\allowbreak 7,\allowbreak 8,\allowbreak 9,\allowbreak 13,\allowbreak 15,\allowbreak 16,\allowbreak 18,\allowbreak 19,\allowbreak 20,\allowbreak 21,\allowbreak 24\}$, there is no residue pair $(a,b)$ modulo $2$ with $F(a,b)\equiv0$ and $H_i(a,b)\equiv0\pmod{2}$, as one checks for the $2^2$ residue pairs; hence $H_i\neq0$. For $i\in\{11,\allowbreak 12,\allowbreak 22,\allowbreak 23\}$, there is no residue pair $(a,b)$ modulo $3$ with $F(a,b)\equiv0$ and $H_i(a,b)\equiv0\pmod{3}$, as one checks for the $3^2$ residue pairs; hence $H_i\neq0$. For $i\in\{14,\allowbreak 17\}$, there is no residue pair $(a,b)$ modulo $4$ with $F(a,b)\equiv0$ and $H_i(a,b)\equiv0\pmod{4}$, as one checks for the $4^2$ residue pairs; hence $H_i\neq0$. For $i\in\{3,\allowbreak 10\}$, there is no residue pair $(a,b)$ modulo $5$ with $F(a,b)\equiv0$ and $H_i(a,b)\equiv0\pmod{5}$, as one checks for the $5^2$ residue pairs; hence $H_i\neq0$. Also, all $c_i$ are non-zero.

\paragraph{Step 2: reciprocity.} For every place $v$ of $K$, put
\begin{equation}\label{eq:B}
B_v=\sum_{(i,j)}w_{ij}\,\beta_v(H_i,H_j)+\sum_i\beta_v(H_i,c_i)\in\Z/3\Z,
\end{equation}
the first sum over the edges, and, for a rational prime $p$, put $B_p=\sum_{v\mid p}B_v$. Applying (H3) to every pair $(H_i,H_j)$ with an edge, multiplied by $w_{ij}$, and to every pair $(H_i,c_i)$, and adding, we find that $B_v=0$ for all but finitely many $v$ and $\sum_pB_p=0$, as the complex place contributes $0$. By (H1) and the definition of $c_p$, $B_p=\sum_{(i,j)}w_{ij}\beta_p(c_p(H_i),c_p(H_j))+\sum_i\beta_p(c_p(H_i),c_p(c_i))$, where $\beta_p$ is the form of Lemma~\ref{lem:M3} if $p=3$.

\paragraph{Step 3: the primes outside $S$.} Let $S=\{2,3\}$. It contains $3$ and the primes dividing the norms of the constants $c_i$, of the $D_i$, of the $R_{ij}$. Let $p\notin S$, and let $v$ be a place above $p$, with residue field $k_v$; then all these numbers are units at $v$. First, no edge $(i,j)$ has both values divisible by $v$. Indeed, suppose that $v$ divides $H_i$ and $H_j$. If $R_{ij}$ is defined, let $k$ be the linear endpoint used for it and $\{k,l\}=\{i,j\}$, and put $t=\xi$. As $H_k=\varepsilon_k(\lambda_k\eta-s_k(\xi))$ and $\lambda_k$ is a unit at $v$, we have $\eta\equiv s_k(t)/\lambda_k$ modulo $v$, hence $f_k(t)\equiv\lambda_k^dF(x,y)=0$ and $g_{kl}(t)\equiv\lambda_k^{d_l}H_l\equiv0$ modulo $v$. The resultant can be written as $R_{ij}=Af_k+Cg_{kl}$ with polynomials $A$ and $C$ with coefficients in $\Z[\omega]$, so $v$ divides $R_{ij}$, a contradiction. By (H2), a term of \eqref{eq:B} in which all arguments are units vanishes. Now let $v$ divide $H_i$. Collecting the terms of \eqref{eq:B} which contain $H_i$, and using (H1) and $\rho_{ji}=-w_{ij}$ for the edges $(j,i)$ with $j<i$, we obtain $\beta_v(H_i,\Pi_i)$ with $\Pi_i=c_i\prod_jH_j^{\rho_{ij}}$, which is a unit. We claim that $D_i^3\Pi_i\equiv W^3$ modulo $v$ for some $W\in K$ integral at $v$. If $i$ is linear, put $t=\xi$; then $g_{ij}(t)\equiv\lambda_i^{d_j}H_j$ and $f_i(t)\equiv0$ modulo $v$ as above, and \eqref{eq:starlin} gives $D_i^3\lambda_i^{E_i}\Pi_i\equiv w_i(t)^3$ modulo $v$, where $E_i=e_i+\sum_j\rho_{ij}d_j$ is divisible by $3$; this gives the claim with $W=w_i(t)/\lambda_i^{E_i/3}$. Otherwise, \eqref{eq:star} gives the claim with $W=W_i(x,y)$. Hence the reduction of $\Pi_i$ is a non-zero cube in $k_v$, $\chi(\Pi_i)=0$, and $\beta_v(H_i,\Pi_i)=v(H_i)\chi(\Pi_i)=0$ by (H2). Hence $B_v=0$, and $B_p=0$.

\paragraph{Step 4: the primes in $S$.} Let $p\in S$. Here $2$ is inert. Suppose that $x\equiv a$ and $y\equiv b\pmod{p^k}$. Then $H_i\equiv H_i(a,b)\pmod{p^k}$, and Lemma~\ref{lem:residue}, applied with $\zeta$ the residue of $H_i(a,b)$, gives a vector $b_i$ and a subspace $D_i$ with $c_p(H_i)\in b_i+D_i$; when $H_i$ has rational coefficients, $H_i\in\Z$, and we use this in (i) and (iv). With $\gamma_i=c_p(c_i)$, if the hypotheses of Lemma~\ref{lem:envelope} hold, then $B_p$ is determined by $a$, $b$ and $k$. Starting from the solutions of \eqref{eq:main} modulo $p$, we refine a residue class $x\equiv a$, $y\equiv b\pmod{p^k}$ to the classes modulo $p^{k+1}$ that contain solutions of \eqref{eq:main} modulo $p^{k+1}$, until these hypotheses hold, or until no solution remains. Table~\ref{tab:local} lists the resulting terminal classes and the value of $B_p$ in each; every integer solution lies in one of them. The result is $B_{2}=0$, $B_{3}=1$.

\paragraph{Step 5: contradiction.} By Steps~2 and~3, $\sum_{p\in S}B_p=0$. But by Step~4, $\sum_{p\in S}B_p=1$ in $\Z/3\Z$. This contradiction proves the theorem.
\qed


\begin{thebibliography}{9}
\bibitem{N} J.~Neukirch, \emph{Algebraic Number Theory}, Grundlehren der mathematischen Wissenschaften 322, Springer, Berlin, 1999.
\end{thebibliography}

\clearpage
\begingroup\small
\setlength{\LTcapwidth}{\textwidth}
\begin{longtable}{@{}rlll@{}}
\caption{The polynomials $H_i$, the constants $c_i$ and, for the linear vertices, $\eta=r_i(\xi)$. Here $\langle A,B\rangle=A+B\omega$.}\label{tab:vert1}\\
\hline $i$ & $H_i$ & $c_i$ & chart\\\hline\endfirsthead\hline $i$ & $H_i$ & $c_i$ & chart\\\hline\endhead\hline\endfoot
0 & $\langle -y+1,1\rangle$ & $\langle 1,-1\rangle$ & $y=\langle 1,1\rangle$\\
1 & $\langle -y,1\rangle$ & $\langle 2,0\rangle$ & $y=\langle 0,1\rangle$\\
2 & $\langle -y,-1\rangle$ & $\langle -3,-3\rangle$ & $y=\langle 0,-1\rangle$\\
3 & $\langle -y-1,0\rangle$ & $\langle -1,-1\rangle$ & $y=\langle -1,0\rangle$\\
4 & $\langle 2x+1,1\rangle$ & $\langle 1,0\rangle$ & --\\
5 & $\langle 2x,1\rangle$ & $\langle -1,-1\rangle$ & --\\
6 & $\langle 2x,-1\rangle$ & $\langle -1,-1\rangle$ & --\\
7 & $\langle 2x-1,-1\rangle$ & $\langle -1,-1\rangle$ & --\\
8 & $\langle 2x+2,1\rangle$ & $\langle 0,-2\rangle$ & --\\
9 & $\langle 2x+1,-1\rangle$ & $\langle 0,-4\rangle$ & --\\
10 & $\langle 2x-y,0\rangle$ & $\langle 0,-1\rangle$ & $y=\langle 2x,0\rangle$\\
11 & $\langle 2x,-y\rangle$ & $\langle 0,-12\rangle$ & $y=\langle -2x,-2x\rangle$\\
12 & $\langle 2x+y,y\rangle$ & $\langle 2,-2\rangle$ & $y=\langle 0,2x\rangle$\\
13 & $\langle 2x+1,-y+1\rangle$ & $\langle -2,-1\rangle$ & $y=\langle -2x,-2x-1\rangle$\\
14 & $\langle 2x+y+1,y+1\rangle$ & $\langle -3,0\rangle$ & $y=\langle -1,2x\rangle$\\
15 & $\langle 2x+y,1\rangle$ & $\langle -1,-1\rangle$ & $y=\langle -2x,-1\rangle$\\
16 & $\langle 2x-y,-1\rangle$ & $\langle -4,-4\rangle$ & $y=\langle 2x,-1\rangle$\\
17 & $\langle 2x,-y-1\rangle$ & $\langle -1,-2\rangle$ & $y=\langle -2x-1,-2x\rangle$\\
18 & $\langle 2x+y,y-1\rangle$ & $\langle 0,-3\rangle$ & $y=\langle 1,2x+1\rangle$\\
19 & $\langle 2x+y-1,-1\rangle$ & $\langle 0,-1\rangle$ & $y=\langle -2x+1,1\rangle$\\
20 & $\langle -y-1,-1\rangle$ & $\langle 4,0\rangle$ & $y=\langle -1,-1\rangle$\\
21 & $\langle 2x-y+1,1\rangle$ & $\langle -2,-2\rangle$ & $y=\langle 2x+1,1\rangle$\\
22 & $\langle 3x+1,2\rangle$ & $\langle -6,0\rangle$ & $3x=\langle -1,-2\rangle$\\
23 & $\langle 3x-1,-2\rangle$ & $\langle -4,-8\rangle$ & $3x=\langle 1,2\rangle$\\
24 & $\langle x,0\rangle$ & $\langle 1,0\rangle$ & $x=\langle 0,0\rangle$\\
\end{longtable}
\endgroup

\begingroup\scriptsize
\setlength{\LTcapwidth}{\textwidth}
\begin{longtable}{@{}ll@{}}
\caption{The data of the identities \eqref{eq:starlin} and \eqref{eq:star}; the entry $j{:}e$ means $\rho_{ij}=e\neq0$.}\label{tab:star}\\
\hline\endfirsthead\hline\endhead\hline\endfoot
\multicolumn{2}{@{}p{0.97\textwidth}@{}}{$i=0$, $D_i=1$, $\rho_{ij}$: $4{:}1, \allowbreak 10{:}1, \allowbreak 11{:}2, \allowbreak 15{:}2, \allowbreak 16{:}1$}\\
$w_i$ & $\langle -8t^{3}-2t-2,$\\
 & $\quad -8t^{3}-4t^{2}-2t-1\rangle$\\
\addlinespace[2pt]
\multicolumn{2}{@{}p{0.97\textwidth}@{}}{$i=1$, $D_i=1$, $\rho_{ij}$: $4{:}1, \allowbreak 5{:}2, \allowbreak 6{:}1, \allowbreak 15{:}2, \allowbreak 19{:}1, \allowbreak 24{:}2$}\\
$w_i$ & $\langle -8t^{3}-2t+2,$\\
 & $\quad -2t\rangle$\\
\addlinespace[2pt]
\multicolumn{2}{@{}p{0.97\textwidth}@{}}{$i=2$, $D_i=1$, $\rho_{ij}$: $6{:}2, \allowbreak 10{:}2, \allowbreak 12{:}1, \allowbreak 19{:}1, \allowbreak 21{:}2$}\\
$w_i$ & $\langle 4t^{2}+4t+3,$\\
 & $\quad -4t^{2}+2t\rangle$\\
\addlinespace[2pt]
\multicolumn{2}{@{}p{0.97\textwidth}@{}}{$i=3$, $D_i=3$, $\rho_{ij}$: $13{:}2, \allowbreak 16{:}1, \allowbreak 18{:}1, \allowbreak 21{:}2$}\\
$w_i$ & $\langle 8t^{3}+8t^{2}-4t+3,$\\
 & $\quad 16t^{3}+16t^{2}+4t+6\rangle$\\
\addlinespace[2pt]
\multicolumn{2}{@{}p{0.97\textwidth}@{}}{$i=4$, $D_i=1$, $\rho_{ij}$: $0{:}2, \allowbreak 1{:}2, \allowbreak 15{:}1, \allowbreak 17{:}2, \allowbreak 18{:}1, \allowbreak 20{:}2, \allowbreak 21{:}2$}\\
$W$ & $\langle -3y+2,$\\
 & $\quad -2y^{2}-y+1\rangle$\\
$A$ & $\langle 8x^{2}y^{6}-8x^{2}y^{4}-32x^{2}y^{3}-24x^{2}y^{2}-16x^{2}y-12xy^{5}-24xy^{4}-36xy^{3}$\\
 & $\quad -24xy^{2}-12xy-5y^{8}-5y^{7}+21y^{6}+27y^{5}-22y^{4}-34y^{3}+5y^{2}-21y+1,$\\
 & $\quad -16x^{2}y^{5}-24x^{2}y^{4}-32x^{2}y^{3}-8x^{2}y^{2}+8x^{2}-4xy^{7}-8xy^{6}-24xy^{5}-24xy^{4}$\\
 & $\quad -24xy^{3}+4xy+8x-y^{9}-4y^{8}+7y^{7}+21y^{6}-12y^{5}-48y^{4}-10y^{3}-2y^{2}-13y+5\rangle$\\
$B$ & $\langle 12x^{3}y^{8}+44x^{3}y^{7}-28x^{3}y^{6}-60x^{3}y^{5}+40x^{3}y^{4}-8x^{3}y^{3}-156x^{3}y^{2}$\\
 & $\quad +12x^{3}y-12x^{3}-2x^{2}y^{9}+18x^{2}y^{8}+40x^{2}y^{7}+48x^{2}y^{6}+54x^{2}y^{5}$\\
 & $\quad -76x^{2}y^{4}-148x^{2}y^{3}-30x^{2}y^{2}+5xy^{9}+10xy^{8}+21xy^{7}+41xy^{6}+14xy^{5}-6xy^{4}$\\
 & $\quad +18xy^{3}+6xy^{2}-21xy-x+y^{11}-5y^{10}-10y^{9}+13y^{8}+44y^{7}+50y^{6}+11y^{5}+22y^{4}$\\
 & $\quad -9y^{3}-34y^{2}+5y-1,$\\
 & $\quad 4x^{3}y^{9}+16x^{3}y^{8}-4x^{3}y^{7}-84x^{3}y^{6}-24x^{3}y^{5}+24x^{3}y^{4}-128x^{3}y^{3}$\\
 & $\quad -64x^{3}y^{2}+52x^{3}y+4x^{3}+4x^{2}y^{9}+26x^{2}y^{8}+10x^{2}y^{7}+6x^{2}y^{6}-70x^{2}y^{5}$\\
 & $\quad -152x^{2}y^{4}-92x^{2}y^{3}+50x^{2}y^{2}-2x^{2}y+10x^{2}+2xy^{10}+5xy^{9}+11xy^{8}+14xy^{7}$\\
 & $\quad +14xy^{6}-31xy^{5}-8xy^{4}+6xy^{3}-25xy^{2}-14xy+2x-4y^{10}+8y^{9}+35y^{8}+39y^{7}+7y^{6}$\\
 & $\quad -19y^{5}-5y^{4}-44y^{3}-10y^{2}+12y+2\rangle$\\
\addlinespace[2pt]
\multicolumn{2}{@{}p{0.97\textwidth}@{}}{$i=5$, $D_i=1$, $\rho_{ij}$: $1{:}1, \allowbreak 11{:}1, \allowbreak 12{:}1, \allowbreak 21{:}2$}\\
$W$ & $\langle y^{2},$\\
 & $\quad y^{2}-y\rangle$\\
$A$ & $\langle y^{3}-2y^{2}-y+1,$\\
 & $\quad -2y^{2}+y+1\rangle$\\
$B$ & $\langle -4x^{3}y^{3}+8x^{3}y^{2}+12x^{3}y+4x^{3}-2x^{2}y+6x^{2}+xy^{3}-2xy^{2}+2x-y^{4}-y^{3}-y^{2}$\\
 & $\quad +2y,$\\
 & $\quad 8x^{3}y^{2}+4x^{3}y-4x^{3}+2x^{2}y^{3}-4x^{2}y^{2}+4x^{2}y+4x^{2}+xy^{3}-2xy^{2}+3xy+x-y^{4}$\\
 & $\quad +y^{2}+2y-1\rangle$\\
\addlinespace[2pt]
\multicolumn{2}{@{}p{0.97\textwidth}@{}}{$i=6$, $D_i=1$, $\rho_{ij}$: $1{:}2, \allowbreak 2{:}1, \allowbreak 10{:}2, \allowbreak 11{:}1, \allowbreak 13{:}2, \allowbreak 14{:}1, \allowbreak 16{:}1, \allowbreak 19{:}1, \allowbreak 20{:}1$}\\
$W$ & $\langle y^{2}+1,$\\
 & $\quad -y^{2}-y+2\rangle$\\
$A$ & $\langle -32x^{4}y^{4}-32x^{4}y^{3}-32x^{4}y^{2}-16x^{3}y^{5}-32x^{3}y^{4}-64x^{3}y^{3}-32x^{3}y^{2}$\\
 & $\quad -16x^{3}y+16x^{3}+24x^{2}y^{6}-8x^{2}y^{4}-64x^{2}y^{3}-40x^{2}y^{2}-32x^{2}y+8xy^{7}$\\
 & $\quad +24xy^{6}+32xy^{5}+4xy^{4}-20xy^{3}-28xy^{2}-8xy+y^{8}+14y^{6}+6y^{5}+29y^{4}+11y^{3}$\\
 & $\quad -10y^{2}-y-2,$\\
 & $\quad -32x^{4}y^{4}-32x^{4}y^{3}-64x^{4}y^{2}-32x^{4}y-32x^{4}+16x^{3}y^{5}-16x^{3}y^{4}$\\
 & $\quad -32x^{3}y^{3}-64x^{3}y^{2}-32x^{3}y-16x^{3}+16x^{2}y^{6}+24x^{2}y^{5}+32x^{2}y^{4}$\\
 & $\quad -8x^{2}y^{3}-32x^{2}y^{2}-40x^{2}y-16x^{2}-4xy^{7}+28xy^{6}+40xy^{5}+56xy^{4}+20xy^{3}$\\
 & $\quad +4xy^{2}-8xy-4x-y^{9}+2y^{8}-3y^{7}+6y^{6}-5y^{5}+37y^{4}+6y^{3}+14y^{2}-y-5\rangle$\\
$B$ & $\langle -12x^{3}y^{8}+80x^{3}y^{7}+24x^{3}y^{6}+192x^{3}y^{5}-28x^{3}y^{4}+52x^{3}y^{3}-8x^{3}y^{2}$\\
 & $\quad -4x^{3}y-8x^{3}-14x^{2}y^{9}-24x^{2}y^{8}-6x^{2}y^{7}+52x^{2}y^{6}+10x^{2}y^{5}$\\
 & $\quad +102x^{2}y^{4}-32x^{2}y^{3}+40x^{2}y^{2}-26x^{2}y+10x^{2}-2xy^{10}-19xy^{9}+7xy^{8}-47xy^{7}$\\
 & $\quad +52xy^{6}-35xy^{5}+56xy^{4}-21xy^{3}+12xy^{2}-14xy+3x+2y^{11}-3y^{10}-7y^{9}-16y^{8}+2y^{7}$\\
 & $\quad -13y^{6}+41y^{5}-21y^{4}+9y^{3}-12y^{2}+6y,$\\
 & $\quad 4x^{3}y^{9}-24x^{3}y^{8}+12x^{3}y^{7}-24x^{3}y^{6}+172x^{3}y^{5}-28x^{3}y^{4}+152x^{3}y^{3}$\\
 & $\quad -40x^{3}y^{2}+36x^{3}y-4x^{3}-6x^{2}y^{9}-54x^{2}y^{8}-2x^{2}y^{7}-68x^{2}y^{6}+58x^{2}y^{5}$\\
 & $\quad +8x^{2}y^{4}+42x^{2}y^{3}+20x^{2}y^{2}-6x^{2}+4xy^{10}-6xy^{9}-7xy^{8}-42xy^{7}+12xy^{6}$\\
 & $\quad -54xy^{5}+45xy^{4}-25xy^{3}+18xy^{2}-19xy+2x+y^{11}+2y^{10}-13y^{9}-4y^{8}-42y^{7}+y^{6}$\\
 & $\quad -5y^{5}+6y^{4}+5y^{3}-3y^{2}+y+1\rangle$\\
\addlinespace[2pt]
\multicolumn{2}{@{}p{0.97\textwidth}@{}}{$i=7$, $D_i=1$, $\rho_{ij}$: $11{:}2, \allowbreak 12{:}2, \allowbreak 16{:}1, \allowbreak 20{:}2$}\\
$W$ & $\langle -y+2,$\\
 & $\quad y^{2}-y+1\rangle$\\
$A$ & $\langle -4xy^{2}+4x+y^{4}-2y^{3}-9y^{2}+3y+2,$\\
 & $\quad -4xy^{2}-8xy+y^{4}+3y^{3}-6y^{2}-3y+5\rangle$\\
$B$ & $\langle -4x^{3}y^{4}+28x^{3}y^{2}-20x^{3}y-8x^{3}-4x^{2}y^{4}+2x^{2}y^{3}-2x^{2}y^{2}-8x^{2}y$\\
 & $\quad +10x^{2}+2xy^{5}-5xy^{4}-9xy^{3}+2xy^{2}+5xy+5x+y^{6}-9y^{4}+2y^{3}+6y^{2}-y+1,$\\
 & $\quad -4x^{3}y^{4}-20x^{3}y^{3}+8x^{3}y^{2}-4x^{3}y-28x^{3}-6x^{2}y^{4}-8x^{2}y^{3}-2x^{2}y^{2}$\\
 & $\quad -18x^{2}y-4x^{2}+2xy^{5}+2xy^{4}-11xy^{3}-5xy^{2}+5x+y^{6}+5y^{5}-3y^{4}-6y^{3}+6y^{2}\rangle$\\
\addlinespace[2pt]
\multicolumn{2}{@{}p{0.97\textwidth}@{}}{$i=8$, $D_i=1$, $\rho_{ij}$: $11{:}2, \allowbreak 14{:}2, \allowbreak 17{:}2, \allowbreak 19{:}1$}\\
$W$ & $\langle 4y-2,$\\
 & $\quad y^{2}+y-6\rangle$\\
$A$ & $\langle -32x^{2}y-16x^{2}-56xy^{2}-48xy-16x-2y^{4}+9y^{3}+27y^{2}-88y-28,$\\
 & $\quad -32x^{3}-80x^{2}y-32x^{2}-88xy^{2}-80xy-24x+16y^{3}-78y^{2}-112y+8\rangle$\\
$B$ & $\langle 8x^{3}y^{4}-244x^{3}y^{3}-428x^{3}y^{2}+160x^{3}y+64x^{3}-128x^{2}y^{4}-164x^{2}y^{3}$\\
 & $\quad +200x^{2}y^{2}-88x^{2}y-96x^{2}-44xy^{5}-30xy^{4}+89xy^{3}-167xy^{2}-16xy+84x-10y^{6}$\\
 & $\quad +26y^{5}+64y^{4}-69y^{3}+109y^{2}+124y-52,$\\
 & $\quad -304x^{3}y^{3}-56x^{3}y^{2}+272x^{3}y-48x^{3}-108x^{2}y^{4}+82x^{2}y^{3}+114x^{2}y^{2}$\\
 & $\quad -192x^{2}y-16x^{2}-28xy^{5}+90xy^{4}+9xy^{3}-77xy^{2}+160xy+52x-4y^{6}+72y^{5}+14y^{4}$\\
 & $\quad +19y^{3}+173y^{2}-40y-84\rangle$\\
\addlinespace[2pt]
\multicolumn{2}{@{}p{0.97\textwidth}@{}}{$i=9$, $D_i=1$, $\rho_{ij}$: $12{:}1, \allowbreak 14{:}1, \allowbreak 15{:}2, \allowbreak 17{:}1$}\\
$W$ & $\langle 2y-2,$\\
 & $\quad -2y^{2}+2y-4\rangle$\\
$A$ & $\langle 8y^{3}+4y^{2}-8y+12,$\\
 & $\quad -16x+24y^{2}-20y+12\rangle$\\
$B$ & $\langle -32x^{3}y^{3}-16x^{3}y^{2}+64x^{3}y+16x^{3}+16x^{2}y^{3}+104x^{2}y^{2}+56x^{2}y+16x^{2}$\\
 & $\quad +32xy^{3}-28xy-20x+12y^{4}+8y^{3}-12y^{2}+20y+4,$\\
 & $\quad -96x^{3}y^{2}-16x^{3}y-16x^{3}-16x^{2}y^{3}+24x^{2}y^{2}+80x^{2}y+40x^{2}+52xy^{2}-12xy-8x$\\
 & $\quad -4y^{4}+40y^{3}-32y^{2}+28y-16\rangle$\\
\addlinespace[2pt]
\multicolumn{2}{@{}p{0.97\textwidth}@{}}{$i=10$, $D_i=2$, $\rho_{ij}$: $0{:}2, \allowbreak 2{:}1, \allowbreak 6{:}1, \allowbreak 20{:}2, \allowbreak 22{:}1, \allowbreak 23{:}2$}\\
$w_i$ & $\langle 8t^{2}+4,$\\
 & $\quad -24t^{3}+4t^{2}-6t+2\rangle$\\
\addlinespace[2pt]
\multicolumn{2}{@{}p{0.97\textwidth}@{}}{$i=11$, $D_i=1$, $\rho_{ij}$: $0{:}1, \allowbreak 5{:}2, \allowbreak 6{:}2, \allowbreak 7{:}1, \allowbreak 8{:}1, \allowbreak 12{:}1$}\\
$w_i$ & $\langle 8t^{2}+4,$\\
 & $\quad -8t^{2}+2\rangle$\\
\addlinespace[2pt]
\multicolumn{2}{@{}p{0.97\textwidth}@{}}{$i=12$, $D_i=1$, $\rho_{ij}$: $2{:}2, \allowbreak 5{:}2, \allowbreak 7{:}1, \allowbreak 9{:}2, \allowbreak 11{:}2, \allowbreak 15{:}1$}\\
$w_i$ & $\langle 16t^{2}+2,$\\
 & $\quad 8t^{2}-2\rangle$\\
\addlinespace[2pt]
\multicolumn{2}{@{}p{0.97\textwidth}@{}}{$i=13$, $D_i=1$, $\rho_{ij}$: $3{:}1, \allowbreak 6{:}1, \allowbreak 14{:}1, \allowbreak 19{:}1, \allowbreak 24{:}1$}\\
$w_i$ & $\langle -2t+1,$\\
 & $\quad -4t-1\rangle$\\
\addlinespace[2pt]
\multicolumn{2}{@{}p{0.97\textwidth}@{}}{$i=14$, $D_i=2$, $\rho_{ij}$: $6{:}2, \allowbreak 8{:}1, \allowbreak 9{:}2, \allowbreak 13{:}2, \allowbreak 16{:}2, \allowbreak 17{:}1, \allowbreak 19{:}2, \allowbreak 22{:}2$}\\
$w_i$ & $\langle -48t^{3}+168t^{2}-36t-72,$\\
 & $\quad -240t^{3}+120t^{2}+108t-36\rangle$\\
\addlinespace[2pt]
\multicolumn{2}{@{}p{0.97\textwidth}@{}}{$i=15$, $D_i=1$, $\rho_{ij}$: $0{:}1, \allowbreak 1{:}1, \allowbreak 4{:}2, \allowbreak 9{:}1, \allowbreak 12{:}2, \allowbreak 16{:}1, \allowbreak 17{:}2, \allowbreak 18{:}1, \allowbreak 20{:}1$}\\
$w_i$ & $\langle 16t^{3}+24t^{2}-28t+4,$\\
 & $\quad -16t^{3}+40t^{2}-12t-4\rangle$\\
\addlinespace[2pt]
\multicolumn{2}{@{}p{0.97\textwidth}@{}}{$i=16$, $D_i=1$, $\rho_{ij}$: $0{:}2, \allowbreak 3{:}2, \allowbreak 6{:}2, \allowbreak 7{:}2, \allowbreak 14{:}1, \allowbreak 15{:}2$}\\
$w_i$ & $\langle -32t^{3}-48t^{2}+8,$\\
 & $\quad -64t^{3}+24t+4\rangle$\\
\addlinespace[2pt]
\multicolumn{2}{@{}p{0.97\textwidth}@{}}{$i=17$, $D_i=1$, $\rho_{ij}$: $4{:}1, \allowbreak 8{:}1, \allowbreak 9{:}2, \allowbreak 14{:}2, \allowbreak 15{:}1, \allowbreak 18{:}1, \allowbreak 21{:}1, \allowbreak 23{:}1$}\\
$w_i$ & $\langle -24t^{3}-36t^{2}-6t,$\\
 & $\quad -24t^{3}-12t^{2}+6t\rangle$\\
\addlinespace[2pt]
\multicolumn{2}{@{}p{0.97\textwidth}@{}}{$i=18$, $D_i=2$, $\rho_{ij}$: $3{:}2, \allowbreak 4{:}2, \allowbreak 15{:}2, \allowbreak 17{:}2, \allowbreak 24{:}2$}\\
$w_i$ & $\langle -24t^{2}+6,$\\
 & $\quad 24t+6\rangle$\\
\addlinespace[2pt]
\multicolumn{2}{@{}p{0.97\textwidth}@{}}{$i=19$, $D_i=1$, $\rho_{ij}$: $1{:}2, \allowbreak 2{:}2, \allowbreak 6{:}2, \allowbreak 8{:}2, \allowbreak 13{:}2, \allowbreak 14{:}1, \allowbreak 20{:}2, \allowbreak 21{:}2$}\\
$w_i$ & $\langle -96t^{3}-144t^{2}+72t,$\\
 & $\quad -160t^{3}+16t^{2}+72t-24\rangle$\\
\addlinespace[2pt]
\multicolumn{2}{@{}p{0.97\textwidth}@{}}{$i=20$, $D_i=1$, $\rho_{ij}$: $4{:}1, \allowbreak 6{:}2, \allowbreak 7{:}1, \allowbreak 10{:}1, \allowbreak 15{:}2, \allowbreak 19{:}1, \allowbreak 24{:}1$}\\
$w_i$ & $\langle -4t,$\\
 & $\quad -8t^{2}\rangle$\\
\addlinespace[2pt]
\multicolumn{2}{@{}p{0.97\textwidth}@{}}{$i=21$, $D_i=1$, $\rho_{ij}$: $2{:}1, \allowbreak 3{:}1, \allowbreak 4{:}1, \allowbreak 5{:}1, \allowbreak 17{:}2, \allowbreak 19{:}1$}\\
$w_i$ & $\langle -16t^{2}-4t+2,$\\
 & $\quad -8t^{2}-8t-2\rangle$\\
\addlinespace[2pt]
\multicolumn{2}{@{}p{0.97\textwidth}@{}}{$i=22$, $D_i=3$, $e_i=0$, $\rho_{ij}$: $10{:}2, \allowbreak 14{:}1$}\\
$w_i$ & $\langle 18t,$\\
 & $\quad -27t^{2}+9t-18\rangle$\\
\addlinespace[2pt]
\multicolumn{2}{@{}p{0.97\textwidth}@{}}{$i=23$, $D_i=3$, $e_i=0$, $\rho_{ij}$: $10{:}1, \allowbreak 17{:}2$}\\
$w_i$ & $\langle 6,$\\
 & $\quad -6\rangle$\\
\addlinespace[2pt]
\multicolumn{2}{@{}p{0.97\textwidth}@{}}{$i=24$, $D_i=1$, $\rho_{ij}$: $1{:}1, \allowbreak 13{:}2, \allowbreak 18{:}1, \allowbreak 20{:}2$}\\
$w_i$ & $\langle t^{2},$\\
 & $\quad -1\rangle$\\
\addlinespace[2pt]
\end{longtable}
\endgroup

\begingroup\small
\setlength{\LTcapwidth}{\textwidth}
\begin{longtable}{@{}lrl@{\qquad}lrl@{}}
\caption{The edges $(i,j)$ with a linear endpoint, the linear endpoint $k$ used, and the norm of $R_{ij}=\operatorname{Res}_t(f_k,g_{kl})$.}\label{tab:sep}\\
\hline $(i,j)$ & $k$ & $N(R_{ij})$ & $(i,j)$ & $k$ & $N(R_{ij})$\\\hline\endfirsthead\hline $(i,j)$ & $k$ & $N(R_{ij})$ & $(i,j)$ & $k$ & $N(R_{ij})$\\\hline\endhead\hline\endfoot
$(0,4)$ & 0 & $2^{10}$ & $(7,11)$ & 11 & $2^{8}\cdot 3$\\
$(0,10)$ & 0 & $2^{8}\cdot 3$ & $(7,12)$ & 12 & $2^{8}\cdot 3^{2}$\\
$(0,11)$ & 0 & $2^{8}\cdot 3$ & $(7,16)$ & 16 & $2^{8}$\\
$(0,15)$ & 0 & $2^{8}\cdot 3^{2}$ & $(7,20)$ & 20 & $2^{8}$\\
$(0,16)$ & 0 & $2^{10}\cdot 3$ & $(8,11)$ & 11 & $2^{12}$\\
$(1,4)$ & 1 & $2^{8}$ & $(8,14)$ & 14 & $2^{8}\cdot 3^{3}$\\
$(1,5)$ & 1 & $2^{8}$ & $(8,17)$ & 17 & $2^{10}\cdot 3^{2}$\\
$(1,6)$ & 1 & $2^{10}$ & $(8,19)$ & 19 & $2^{14}\cdot 3^{2}$\\
$(1,15)$ & 1 & $2^{14}$ & $(9,12)$ & 12 & $2^{12}$\\
$(1,19)$ & 1 & $2^{8}$ & $(9,14)$ & 14 & $2^{10}\cdot 3^{2}$\\
$(1,24)$ & 1 & $1$ & $(9,15)$ & 15 & $2^{14}\cdot 3^{2}$\\
$(2,6)$ & 2 & $2^{10}$ & $(9,17)$ & 17 & $2^{8}\cdot 3^{3}$\\
$(2,10)$ & 2 & $2^{8}\cdot 3$ & $(10,20)$ & 10 & $2^{10}$\\
$(2,12)$ & 2 & $2^{8}\cdot 3$ & $(10,22)$ & 10 & $2^{2}\cdot 3^{4}$\\
$(2,19)$ & 2 & $2^{8}\cdot 3^{2}$ & $(10,23)$ & 10 & $2^{2}\cdot 3^{4}$\\
$(2,21)$ & 2 & $2^{10}\cdot 3$ & $(11,12)$ & 11 & $2^{8}\cdot 3^{4}$\\
$(3,13)$ & 3 & $2^{8}\cdot 3$ & $(12,15)$ & 12 & $2^{8}$\\
$(3,16)$ & 3 & $2^{12}\cdot 3$ & $(13,14)$ & 13 & $2^{8}\cdot 3^{5}$\\
$(3,18)$ & 3 & $2^{8}\cdot 3$ & $(13,19)$ & 13 & $2^{12}\cdot 3$\\
$(3,21)$ & 3 & $2^{12}\cdot 3$ & $(13,24)$ & 13 & $3$\\
$(4,15)$ & 15 & $2^{8}$ & $(14,16)$ & 14 & $2^{8}\cdot 3^{4}$\\
$(4,17)$ & 17 & $2^{10}\cdot 3^{2}$ & $(14,17)$ & 14 & $2^{8}\cdot 3^{6}$\\
$(4,18)$ & 18 & $2^{8}\cdot 3^{3}$ & $(14,19)$ & 14 & $2^{8}\cdot 3^{3}$\\
$(4,20)$ & 20 & $2^{10}$ & $(14,22)$ & 14 & $2^{2}\cdot 3^{4}$\\
$(4,21)$ & 21 & $2^{8}\cdot 3$ & $(15,16)$ & 15 & $2^{16}\cdot 3$\\
$(5,11)$ & 11 & $2^{8}\cdot 3^{2}$ & $(15,17)$ & 15 & $2^{8}\cdot 3^{3}$\\
$(5,12)$ & 12 & $2^{8}\cdot 3$ & $(15,18)$ & 15 & $2^{12}\cdot 3$\\
$(5,21)$ & 21 & $2^{8}$ & $(15,20)$ & 15 & $2^{8}$\\
$(6,10)$ & 10 & $2^{10}$ & $(17,18)$ & 17 & $2^{8}\cdot 3^{5}$\\
$(6,11)$ & 11 & $2^{8}$ & $(17,21)$ & 17 & $2^{8}\cdot 3^{4}$\\
$(6,13)$ & 13 & $2^{8}\cdot 3^{3}$ & $(17,23)$ & 17 & $2^{2}\cdot 3^{4}$\\
$(6,14)$ & 14 & $2^{10}\cdot 3^{2}$ & $(18,24)$ & 18 & $3$\\
$(6,16)$ & 16 & $2^{8}\cdot 3$ & $(19,20)$ & 19 & $2^{14}$\\
$(6,19)$ & 19 & $2^{8}$ & $(19,21)$ & 19 & $2^{16}\cdot 3$\\
$(6,20)$ & 20 & $2^{8}$ & $(20,24)$ & 20 & $1$\\
\end{longtable}
\endgroup

\begingroup\small
\setlength{\LTcapwidth}{\textwidth}
\begin{longtable}{@{}rrl>{\raggedright\arraybackslash}p{0.66\textwidth}@{}}
\caption{The terminal residue classes $x\equiv a$, $y\equiv b\pmod{p^k}$ of Step~4, grouped by $p$, $k$ and the value of $B_p$.}\label{tab:local}\\
\hline $p$ & $k$ & $B_p$ & classes $(a,b)$\\\hline\endfirsthead
\hline $p$ & $k$ & $B_p$ & classes $(a,b)$\\\hline\endhead\hline\endfoot
2 & 2 & $0$ & $(1,0)$,\allowbreak\ $(3,1)$,\allowbreak\ $(3,2)$,\allowbreak\ $(3,3)$\\
\addlinespace[3pt]
3 & 3 & $1$ & $(0,20)$,\allowbreak\ $(1,11)$,\allowbreak\ $(2,6)$,\allowbreak\ $(3,26)$,\allowbreak\ $(4,20)$,\allowbreak\ $(5,18)$,\allowbreak\ $(6,5)$,\allowbreak\ $(7,2)$,\allowbreak\ $(8,3)$,\allowbreak\ $(9,11)$,\allowbreak\ $(10,11)$,\allowbreak\ $(11,15)$,\allowbreak\ $(12,17)$,\allowbreak\ $(13,20)$,\allowbreak\ $(14,0)$,\allowbreak\ $(15,23)$,\allowbreak\ $(16,2)$,\allowbreak\ $(17,12)$,\allowbreak\ $(18,2)$,\allowbreak\ $(19,11)$,\allowbreak\ $(20,24)$,\allowbreak\ $(21,8)$,\allowbreak\ $(22,20)$,\allowbreak\ $(23,9)$,\allowbreak\ $(24,14)$,\allowbreak\ $(25,2)$,\allowbreak\ $(26,21)$\\
\addlinespace[3pt]
\end{longtable}
\endgroup
\end{document}
